% Title text and style
\title{\titlestyle{Microclimate and climatic context determine the phenology of individual trees within cities}}

% Author markers: * corresponding author, dagger equal contribution.
\author[1,2,$\ast$,$\dagger$]{Jiali~Zhu}
\author[2,$\dagger$]{A.~Scientist}
\author[2]{Someone~E.~Else}
\affil[1]{\small Department, Institution, City \& Postal Code, Country.}
\affil[2]{\small Another Department, Different Institution, City \& Postal Code, Country.}
\affil[$\ast$]{\small Corresponding author: example@mail.com}
\affil[$\dagger$]{\small These authors contributed equally to this work.}
\date{}

\maketitle
\thispagestyle{empty}

% Push front-page metadata to the bottom of page 1.
\vfill

\begin{center}
  {\sc Number of words: } abstract XX, main text XX\\
  {\sc Number of references: } \total{citnum}\\
  {\sc Number of figures: } XX colored figures
\end{center}

% Start abstract on a fresh page after the title page.
\clearpage

\section*{Abstract}

The unprecedented urbanization is affecting urban trees and their seasonal life cycles, known as phenology. The diverse species composition and human-driven heterogeneous microclimate further generate pronounced within-city phenological variations. However, such fine-scale seasonality remains poorly studied, constraining management for localized ecosystem services and disservices, such as cooling and pollen control. To address this, we used high-resolution remote sensing products from PlanetScope to infer the spring and fall timing of individual trees documented in urban inventories, and evaluated the species-level response to microclimate conditions characterized by urban high-resolution air temperature data. We further analyzed how these responses varied with the mismatch between regional climatic conditions and species' climatic origins across the contiguous U.S. Our results demonstrated that intra-urban phenological variation rivals inter-city differences, with 51.5\% of species-city combinations exhibiting advanced spring, and 66.5\% exhibiting delayed fall under warmer microclimate. Moreover, the direction and magnitude of microclimate effects are jointly constrained by regional thermal and moisture context and species' climatic origins, leading to divergent responses across cities and taxa. By uncovering the ecological mechanisms underlying intra-urban seasonality heterogeneity and species-specific responses, our findings could enhance insights into urban forestry management and species selection, thereby supporting human well-being and the development of resilient urban forests.

% {\sc Keywords}: keyword1, keyword2, keyword3

% Start main text on a new page.
\clearpage
